\section{Introduction}

Robotics is an important section of engineering which focuses on the design, development and control of machines capable of performing tasks autonomously. Robots are commonly being used in modern applications that can replace humans in a number of occupations by their own structures and operations. Especially industrial robotic designs, they combine mechanical structures, electronic components, sensors, actuators and other computational techniques and are used for the industrial tasks like transporting and unloading materials, assembling, measuring, etc. Therefore, the demand for this kind of robotic designs is gradually increasing and manufacturing of the types of robots are increasing with higher precision and greater productivity. 

Robotic arms are one of the robotic systems that are widely used for material handling, assembly, manufacturing and pick-and-place operations in various industries. It functions more similarly to a human arm. A robotic arm consists of connected links, joints and gripper mechanism which provide movement and pick-and-place operations within the defined workspace.  In robotics, Degree Of Freedom (DOF) is defined as the number of independent axes that a robot can move. In most of the manufacturing industry, robotic arms have replaced human hands with applications like assembling, repairing and pick-and-place goods according to their assigned placements. 
As mentioned before, pick-and-place is one of the fundamental applications of robotic arms in industry where an object is picked from one location and transferred to another. All joints of the robotic arm should be accurately positioned for this task to succeed. The accurate joint configurations and the coordinate system is calculated by using inverse kinematics. Inverse kinematics establishes the relationship between the desired position of the end effector and the corresponding joint angles of the robotic arm. 

The purpose of this project is focused on the design and implementation of the 5 DOF robotic arm with the goal of performing an automated pick-and-place operation. This system integrates mechanical structure which designed using SolidWorks, servo motors, ESP32 microcontroller and a gripper mechanism Considering the structure of this paper, literature review, system overview and architecture, mechanical design and prototype development, Inverse kinematics model and angle calculation , electronic system and microcontroller implementation, experimental results and discussion and finally future improvements and the references. 
